\documentclass[12pt,letterpaper]{article}

% ============================================================
% TwinSight X500 - Informe Final
% Archivo maestro DEFINITIVO - APA 7 UNAD (Instructivo marzo 2023)
% ============================================================

% --- PAQUETES ---
\usepackage[utf8]{inputenc}
\usepackage[T1]{fontenc}
\usepackage[spanish]{babel}
\usepackage{times}
\usepackage[margin=2.54cm]{geometry}   % Margenes 2.54 cm (p. 8)
\usepackage{setspace}
\usepackage{graphicx}
\usepackage{booktabs}
\usepackage{array}
\usepackage{longtable}
\usepackage{float}
\usepackage{placeins}                  % \FloatBarrier SOLO para casos puntuales
\usepackage{caption}
\usepackage{amsmath}
\usepackage{amsfonts}
\usepackage{amssymb}
\usepackage{xurl}
\usepackage[hidelinks]{hyperref}
\usepackage{ragged2e}
\usepackage{titlesec}
\usepackage{tocloft}
\usepackage{enumitem}
\usepackage{etoolbox}

% Interlineado reducido permitido dentro de tablas (p. 23)
\AtBeginEnvironment{table}{\setstretch{1.15}}
\AtBeginEnvironment{longtable}{\setstretch{1.15}}

% --- LISTAS: vineta alineada con la sangria de parrafo (1.27 cm) ---
\setlist[itemize]{leftmargin=1.9cm, itemindent=0pt, labelsep=0.63cm, label=\textbullet, itemsep=0pt, parsep=0pt, topsep=0pt}
\setlist[enumerate]{leftmargin=1.9cm, itemindent=0pt, labelsep=0.63cm, itemsep=0pt, parsep=0pt, topsep=0pt}

% --- ROTULOS EN ESPANOL ---
\addto\captionsspanish{\renewcommand{\tablename}{Tabla}}
\addto\captionsspanish{\renewcommand{\figurename}{Figura}}
\addto\captionsspanish{\renewcommand{\contentsname}{Tabla de Contenido}}
\addto\captionsspanish{\renewcommand{\listtablename}{Lista de Tablas}}
\addto\captionsspanish{\renewcommand{\listfigurename}{Lista de Figuras}}

% --- ESPACIADO GLOBAL (pp. 7-8) ---
\setstretch{2}                          % Interlineado 2.0 en TODO el documento
\setlength{\parindent}{1.27cm}          % Sangria de primera linea
\setlength{\parskip}{0pt}               % Sin espacio libre entre parrafos (p. 7)
\setlength{\emergencystretch}{3em}

\setcounter{secnumdepth}{0}             % Titulos sin numeracion visible (p. 8)
\setcounter{tocdepth}{3}                % TOC hasta Nivel 3 (como el TOC del propio Instructivo)

\hypersetup{bookmarksopen=true, bookmarksopenlevel=1, bookmarksdepth=3}

% Pagina nueva EXCLUSIVAMENTE para titulos de Nivel 1 (p. 8)
\pretocmd{\section}{\clearpage}{}{}
\pretocmd{\subsection}{\phantomsection}{}{}
\pretocmd{\subsubsection}{\phantomsection}{}{}
\pretocmd{\paragraph}{\phantomsection}{}{}
\pretocmd{\subparagraph}{\phantomsection}{}{}
\makeatletter
\providecommand{\ttl@finishall}{}
\providecommand{\contentsfinish}{}
\makeatother

% --- PAGINACION: arabiga, superior derecha, desde la portada (p. 7) ---
\usepackage{fancyhdr}
\pagestyle{fancy}
\fancyhf{}
\renewcommand{\headrulewidth}{0pt}
\setlength{\headheight}{14.5pt}
\fancyhead[R]{\thepage}
\fancypagestyle{plain}{%
  \fancyhf{}%
  \renewcommand{\headrulewidth}{0pt}%
  \fancyhead[R]{\thepage}%
}

% --- FORMATO TOC / LISTA DE TABLAS / LISTA DE FIGURAS ---
\renewcommand{\cftsecpresnum}{\relax}
\renewcommand{\cftsecnumwidth}{0pt}
\renewcommand{\cftsubsecpresnum}{\relax}
\renewcommand{\cftsubsecnumwidth}{0pt}
\renewcommand{\cftsubsubsecpresnum}{\relax}
\renewcommand{\cftsubsubsecnumwidth}{0pt}
\renewcommand{\cftsecleader}{\cftdotfill{\cftdotsep}}
\renewcommand{\cftsubsecleader}{\cftdotfill{\cftdotsep}}
\renewcommand{\cftsubsubsecleader}{\cftdotfill{\cftdotsep}}

\RaggedRight                            % Alineado a la izquierda, sin justificar (p. 8)

\setlength{\cftbeforetoctitleskip}{0pt}
\setlength{\cftaftertoctitleskip}{0pt}
\renewcommand{\cfttoctitlefont}{\hfill\normalfont\bfseries\normalsize}
\renewcommand{\cftaftertoctitle}{\null\hfill}

\setlength{\cftbeforelottitleskip}{0pt}
\setlength{\cftafterlottitleskip}{0pt}
\renewcommand{\cftlottitlefont}{\hfill\normalfont\bfseries\normalsize}
\renewcommand{\cftafterlottitle}{\null\hfill}

\setlength{\cftbeforeloftitleskip}{0pt}
\setlength{\cftafterloftitleskip}{0pt}
\renewcommand{\cftloftitlefont}{\hfill\normalfont\bfseries\normalsize}
\renewcommand{\cftafterloftitle}{\null\hfill}

\renewcommand{\cftsecfont}{\normalfont}
\renewcommand{\cftsecpagefont}{\normalfont}
\renewcommand{\cftsubsecfont}{\normalfont}
\renewcommand{\cftsubsecpagefont}{\normalfont}
\renewcommand{\cftsubsubsecfont}{\normalfont}
\renewcommand{\cftsubsubsecpagefont}{\normalfont}

\setlength{\cftsecindent}{0em}
\setlength{\cftsubsecindent}{1.5em}
\setlength{\cftsubsubsecindent}{3em}

\newlength{\mytocskip}
\setlength{\mytocskip}{0pt}
\setlength{\cftbeforesecskip}{\mytocskip}
\setlength{\cftbeforesubsecskip}{\mytocskip}
\setlength{\cftbeforesubsubsecskip}{\mytocskip}

% Ancho compacto del rotulo: cabe "Figura 10" / "Tabla 10" (3.9em aprox.)
% sin dejar el hueco que marco el revisor. NO bajar de 4.5em.
\renewcommand{\cfttabpresnum}{\tablename~}
\renewcommand{\cfttabfont}{\normalfont}
\renewcommand{\cfttabpagefont}{\normalfont}
\renewcommand{\cfttableader}{\cftdotfill{\cftdotsep}}
\setlength{\cfttabnumwidth}{4.5em}
\setlength{\cftbeforetabskip}{0.15\baselineskip}

\renewcommand{\cftfigpresnum}{\figurename~}
\renewcommand{\cftfigfont}{\normalfont}
\renewcommand{\cftfigpagefont}{\normalfont}
\renewcommand{\cftfigleader}{\cftdotfill{\cftdotsep}}
\setlength{\cftfignumwidth}{4.5em}
\setlength{\cftbeforefigskip}{0.15\baselineskip}

% --- TITULOS APA 7 UNAD (p. 21) ---
% REGLA DE ORO: sin espacio vertical extra alrededor de titulos.
% El titulo ya ocupa una linea a interlineado 2.0; cualquier valor
% adicional aqui es un "espacio libre" prohibido (p. 7) y es lo que
% el revisor ordeno eliminar.
% La forma ESTRELLADA de \titlespacing suprime automaticamente la
% sangria del primer parrafo tras el titulo (regla "sangria a partir
% del segundo parrafo", pp. 7, 18, 19, 25, 26). NO usar \noindent
% manuales en los capitulos.

% Nivel 1: centrado, negrita, mayusculas capitalizadas, pagina nueva
\titleformat{\section}{\normalfont\bfseries\centering}{}{0em}{}
\titlespacing*{\section}{0pt}{0pt}{0pt}

% Nivel 2: izquierda, negrita, mayusculas capitalizadas
\titleformat{\subsection}{\normalfont\bfseries}{}{0em}{}
\titlespacing*{\subsection}{0pt}{0pt}{0pt}

% Nivel 3: izquierda, negrita, cursiva, mayusculas capitalizadas
\titleformat{\subsubsection}{\normalfont\bfseries\itshape}{}{0em}{}
\titlespacing*{\subsubsection}{0pt}{0pt}{0pt}

% Nivel 4: sangria 1.27 cm, negrita, punto final, texto en la misma linea.
% Separacion tras el punto = un espacio normal (0.5em), NO un cuadratin (1em).
\titleformat{\paragraph}[runin]{\normalfont\bfseries}{}{0em}{}[.]
\titlespacing*{\paragraph}{1.27cm}{0pt}{0.5em}

% Nivel 5: igual que Nivel 4, con cursiva
\titleformat{\subparagraph}[runin]{\normalfont\bfseries\itshape}{}{0em}{}[.]
\titlespacing*{\subparagraph}{1.27cm}{0pt}{0.5em}

% --- CAPTIONS Y NOTAS (pp. 22-24) ---
% Linea 1: "Tabla 1" / "Figura 1" en negrita, alineado a la izquierda.
% Linea 2: titulo en cursiva, mayusculas capitalizadas, SIN punto final.
\DeclareCaptionLabelFormat{table_num_bold_left}{\textbf{#1~#2}}
\DeclareCaptionLabelSeparator{mediumspace}{\par\vspace{0.25\baselineskip}}

\captionsetup[table]{
    labelformat=table_num_bold_left,
    labelsep=mediumspace,
    textfont=it,
    justification=raggedright,
    singlelinecheck=false,
    skip=0pt,
    position=top
}
\captionsetup[figure]{
    labelformat=table_num_bold_left,
    labelsep=mediumspace,
    textfont=it,
    justification=raggedright,
    singlelinecheck=false,
    skip=0pt,
    position=top
}

% Nota de tabla/figura: pegada al objeto, interlineado reducido (permitido, p. 23).
% NUNCA anteponer \vspace manual en los capitulos: este macro ya trae el suyo.
\newcommand{\apanote}[1]{%
    \par\vspace{0.3\baselineskip}%
    \begingroup%
    \setstretch{1.15}%
    \RaggedRight%
    \noindent\textit{Nota}. #1\par%
    \endgroup%
}

\begin{document}

% --- PORTADA (p. 9) ---
% Sin entorno titlepage: la portada se numera como pagina 1 (p. 7).
% Titulo en minusculas y negrita; mayuscula solo en siglas y nombres propios.
\thispagestyle{fancy}
\begin{center}
    \setstretch{2}
    \vspace*{2\baselineskip}
    {\normalfont\bfseries TwinSight X500: prototipo web 3D para visualización técnica de hardware complejo\par}
    \vfill
    {\normalfont Alexander Woodcock Salomón\par}
    \vspace{2\baselineskip}
    {\normalfont Asesor\par
    Gustavo Enrique Vejarano Matiz\par}
    \vfill
    {\normalfont Universidad Nacional Abierta y a Distancia UNAD\par
    Escuela de Ciencias Básicas, Tecnología e Ingeniería ECBTI\par
    Ingeniería Multimedia\par
    2026\par}
\end{center}

% --- RESUMEN (p. 12): un solo parrafo, sin sangria de primera linea ---
\section{Resumen}
\setlength{\parindent}{0pt}
La documentación de hardware técnico complejo, como drones y sistemas electromecánicos multicomponente, presenta limitaciones de comunicación espacial cuando depende de imágenes 2D, manuales PDF y fuentes técnicas fragmentadas. Este informe desarrolla TwinSight X500, un prototipo web 3D interactivo orientado a la visualización técnica, inspección y análisis visual del dron Holybro X500~V2. El trabajo se enmarca en Design Science Research, prototipado iterativo, Teoría de la Carga Cognitiva, interacción 3D y renderizado físicamente basado. La build integra selección jerárquica de piezas, ficha técnica inferior, modos \textit{Inspect}, \textit{Analyze} y \textit{Studio}, corte transversal, vista explosionada, filtros y visualización térmica heurística. El pipeline redujo el modelo desde rutas CAD superiores a 6,5 millones de triángulos hasta 95\,617 triángulos optimizados. La validación combinó profiler WebGL, desempeño en tareas cronometradas, SUS, NASA-TLX Raw y \textit{Think-Aloud} con 12 participantes: SUS promedio de 91,88, NASA-TLX Raw de 8,69 para el visor 3D frente a 19,89 para el soporte 2D, menor tiempo medio en las tres tareas cronometradas y 96 registros tarea-condición completados. El prototipo alcanzó rendimiento estable en escritorio, comportamiento funcional en móviles reales y alcance defendible como \textit{visual product twin}, no como gemelo digital operacional.

% "Palabras clave" con sangria, cursiva y negrita (p. 12). La lista NO lleva punto final.
\hspace{1.27cm}\textit{\textbf{Palabras clave:}} WebGL, Unity, WebAssembly, visualización, usabilidad

% --- ABSTRACT (p. 13) ---
\section{Abstract}
Complex hardware documentation, including drones and multicomponent electromechanical systems, presents spatial communication limits when it depends on 2D images, PDF manuals, and fragmented technical sources. This report develops TwinSight X500, an interactive 3D web prototype for technical visualization, inspection, and visual analysis of the Holybro X500~V2. The work is framed through Design Science Research, iterative prototyping, cognitive load theory, 3D interaction, and physically based rendering. The build includes hierarchical part selection, a bottom technical sheet, \textit{Inspect}, \textit{Analyze}, and \textit{Studio} modes, cross-section, exploded view, category filters, and component-based heuristic thermal visualization. The pipeline reduced the model from CAD routes above 6.5 million triangles to 95,617 optimized triangles. Validation combined WebGL profiling, timed task performance, SUS, NASA-TLX Raw, and Think-Aloud with 12 participants: mean SUS was 91.88, NASA-TLX Raw was 8.69 for the 3D viewer versus 19.89 for the 2D reference, lower mean time in the three timed tasks for the 3D condition, and 96 completed task-condition records. The prototype achieved stable desktop performance, functional behavior on real mobile devices, and a defensible scope as a visual product twin, not as an operational digital twin.

\hspace{1.27cm}\textit{\textbf{Keywords:}} WebGL, Unity, WebAssembly, visualization, usability

\setlength{\parindent}{1.27cm}

% --- LISTAS PRELIMINARES: ORDEN ESTRICTO UNAD (pp. 3, 14-17) ---
% 1. Tabla de Contenido  2. Lista de Tablas  3. Lista de Figuras  4. Lista de Apendices
\clearpage
\tableofcontents

\clearpage
\phantomsection
\addcontentsline{toc}{section}{Lista de Tablas}
\listoftables

\clearpage
\phantomsection
\addcontentsline{toc}{section}{Lista de Figuras}
\listoffigures

\section{Lista de Apéndices}
\noindent Apéndice A \quad Repositorio de Código Fuente \cftdotfill{\cftdotsep} \pageref{apendice:A}\par\vspace{0.15\baselineskip}
\noindent Apéndice B \quad Manual de Usuario \cftdotfill{\cftdotsep} \pageref{apendice:B}\par\vspace{0.15\baselineskip}
\noindent Apéndice C \quad Manual Técnico \cftdotfill{\cftdotsep} \pageref{apendice:C}\par\vspace{0.15\baselineskip}
\noindent Apéndice D \quad Documentación Técnica Activa \cftdotfill{\cftdotsep} \pageref{apendice:D}\par\vspace{0.15\baselineskip}
\noindent Apéndice E \quad Matriz de Desconexiones \cftdotfill{\cftdotsep} \pageref{apendice:E}\par\vspace{0.15\baselineskip}
\noindent Apéndice F \quad Catálogo Canónico del Holybro X500 V2 \cftdotfill{\cftdotsep} \pageref{apendice:F}\par\vspace{0.15\baselineskip}
\noindent Apéndice G \quad Validación Funcional de la Escena Final \cftdotfill{\cftdotsep} \pageref{apendice:G}\par\vspace{0.15\baselineskip}
\noindent Apéndice H \quad Instrumentos y Registros de Evaluación \cftdotfill{\cftdotsep} \pageref{apendice:H}\par\vspace{0.15\baselineskip}
\noindent Apéndice I \quad Estado del Pipeline de Activos \cftdotfill{\cftdotsep} \pageref{apendice:I}\par\vspace{0.15\baselineskip}
\noindent Apéndice J \quad Demo Público \cftdotfill{\cftdotsep} \pageref{apendice:J}\par\vspace{0.15\baselineskip}
\noindent Apéndice K \quad Audio y Alcance No Integrado \cftdotfill{\cftdotsep} \pageref{apendice:K}\par\vspace{0.15\baselineskip}
\noindent Apéndice L \quad Auditoría Matemática y Trazabilidad con Wolfram \cftdotfill{\cftdotsep} \pageref{apendice:L}\par\vspace{0.15\baselineskip}
\noindent Apéndice M \quad Correlación entre la Propuesta y el Informe Final \cftdotfill{\cftdotsep} \pageref{apendice:M}\par\vspace{0.15\baselineskip}
\noindent Apéndice N \quad Resumen de Artefactos Activos del Cierre \cftdotfill{\cftdotsep} \pageref{apendice:N}\par\vspace{0.15\baselineskip}

% --- INFORMACION GENERAL DEL TRABAJO DE GRADO ---
\section{Información General del Trabajo de Grado}

\subsection{Integrantes del Proyecto}

\textit{Nombre del estudiante:} Alexander Woodcock Salomón

\textit{Programa Académico:} Ingeniería Multimedia

\textit{Correo Electrónico:} awoodcocks@unadvirtual.edu.co

\textit{Teléfono:} +57 301 234 5678

\textit{Centro Seccional o Zonal:} ZSUR - Pasto

\textit{Nombre del Asesor:} Gustavo Enrique Vejarano Matiz

\subsection{Datos Específicos del Proyecto}

\textit{Duración del proyecto:} Un (1) año

\textit{Línea de Investigación:} Ingeniería Multimedia, Visualización Técnica 3D Web

\textit{Escuela:} Escuela de Ciencias Básicas, Tecnología e Ingeniería (ECBTI)

\textit{Descriptores palabras clave:} WebGL, Unity, WebAssembly, visualización técnica 3D, optimización, carga de trabajo percibida, HCI, Technical Art, hardware complejo, visual product twin, digital shadow.

% --- CAPITULOS ---
\input{chapters/01_introduccion}
\input{chapters/02_marco_referencia}
% =============================================================================
% Capítulo 3 — Marco Metodológico
% Conforme a APA 7 UNAD (Instructivo marzo 2023)
% =============================================================================

\section{Marco Metodológico}

Este capítulo describe el diseño metodológico adoptado para el desarrollo y la evaluación del prototipo web del Holybro X500~V2. La investigación se formuló como aplicada, con enfoque mixto y predominio cualitativo-formativo. El trabajo articuló diseño, construcción y evaluación de un artefacto digital, por lo que adoptó \textit{Design Science Research} (DSR) como marco principal, complementado con prototipado iterativo, observación de uso, cuestionarios estandarizados y medición de KPIs técnicos. La identificación del problema no se restringió al caso del dron, sino que partió de una necesidad industrial más amplia: la visualización, documentación y preparación semántica de hardware complejo.

\subsection{Tipo de Investigación y Enfoque}

\paragraph{Tipo de Investigación} Aplicada, porque busca resolver una necesidad concreta de visualización técnica en entorno web.

\paragraph{Enfoque} Mixto, con predominio cualitativo-formativo y un componente cuantitativo descriptivo y exploratorio.

\paragraph{Marco Metodológico} DSR para construir y evaluar el artefacto, complementado con observación cualitativa de uso, cuestionarios estandarizados y medición de KPIs técnicos.

\paragraph{Unidad de Análisis} La interacción de usuarios con un prototipo web de visualización técnica del dron Holybro X500~V2, contrastada con un soporte 2D de referencia para tareas equivalentes.

\subsection{Marco de Design Science Research}

La investigación se enmarcó en la epistemología de \textit{Design Science Research}. En términos procedimentales, el proyecto adoptó el modelo de seis fases del DSRM propuesto por Peffers et al. (2007), mientras que la evaluación de rigor y relevancia del artefacto se fundamentó en las siete directrices establecidas por Hevner et al. (2004). De este modo, el proceso de construcción y evaluación del artefacto quedó diferenciado de las pautas con las que se juzga su consistencia metodológica y su aporte disciplinar. La Figura \ref{fig:dsrm_aplicado_proyecto} sintetiza esa aplicación del ciclo DSRM al caso del visor WebGL.

\paragraph{Fase 1: Identificación del Problema y Motivación} La documentación 2D tradicional dificulta la comprensión espacial y funcional de sistemas de hardware con múltiples componentes; además, la industria suele mantener separados CAD, manuales, listas de materiales, procedimientos, sensores y conocimiento experto.

\paragraph{Fase 2: Definición de Objetivos de la Solución} Desarrollar un visor web interactivo orientado a facilitar la exploración y lectura técnica del dron Holybro X500~V2 bajo restricciones reales de rendimiento, entendiendo el prototipo como una capa visual-semántica previa a posibles fases de \textit{digital shadow}.

\paragraph{Fase 3: Diseño y Desarrollo} Construir el prototipo en Unity con un \textit{pipeline} CAD a web, UI responsiva y herramientas de análisis visual.

\paragraph{Fase 4: Demostración} Publicar el prototipo en entorno web y ejecutar tareas representativas de exploración, identificación y análisis.

\paragraph{Fase 5: Evaluación} Analizar de forma formativa el rendimiento técnico, la usabilidad percibida y la carga de trabajo percibida del prototipo.

\paragraph{Fase 6: Comunicación} Documentar el proceso, la arquitectura, las decisiones de diseño y los resultados obtenidos.

\begin{figure}[htbp]
    \caption{Aplicación del Ciclo Design Science Research al Proyecto WebGL}
    \label{fig:dsrm_aplicado_proyecto}
    \includegraphics[width=\textwidth,height=0.45\textheight,keepaspectratio]{figures/chapter3/fig_3_dsrm_aplicado_proyecto.pdf}
    \apanote[6.1pt]{La evaluación retroalimenta ajustes del artefacto, pero no cambia el alcance metodológico: se evalúa un \textit{visual product twin}, no un gemelo digital operacional. Elaboración propia.}
\end{figure}

\subsection{Delimitación Metodológica como Visual Product Twin}

El artefacto evaluado en este informe se delimita como \textit{visual product twin} o modelo visual-semántico de producto, no como \textit{digital twin} operacional. Esta delimitación metodológica evita evaluar el sistema con criterios que no corresponden a su alcance actual. Por tanto, la validación se concentra en viabilidad web, comprensión espacial, inspección visual, usabilidad percibida, carga de trabajo subjetiva y coherencia funcional de la interfaz. No se evalúan sincronización con activo físico, mantenimiento predictivo, integración PLM/CMMS/SCADA ni modelos calibrados de simulación, porque esas capacidades pertenecen al roadmap futuro.

Esta decisión también orienta los instrumentos. Los KPIs técnicos y las pruebas de usuario permiten juzgar si la capa visual-semántica es funcional, comprensible y defendible como base de trabajo; no permiten afirmar que el sistema sea una sombra digital o un gemelo digital. Para evolucionar a esas categorías serían necesarios, como mínimo, un \textit{twin manifest} por pieza, telemetría histórica o en vivo, trazabilidad de eventos, reglas de mantenimiento y validación de datos operacionales (véase la Figura \ref{fig:frontera_metodologica_visual_twin}).

\begin{figure}[htbp]
    \caption{Frontera Metodológica entre lo Evaluado y lo Proyectado}
    \label{fig:frontera_metodologica_visual_twin}
    \includegraphics[width=0.78\textwidth,height=0.45\textheight,keepaspectratio]{figures/chapter3/fig_3_frontera_metodologica.pdf}
    \apanote[-18.3pt]{Las capacidades de telemetría, sincronización, mantenimiento predictivo y simulación calibrada se mantienen fuera de la evaluación actual. Elaboración propia.}
\end{figure}

\subsection{Fases de Trabajo}

\paragraph{Fase I: Análisis y Fundamentación} Revisión bibliográfica sobre WebGL, visualización 3D, teoría de la carga cognitiva, carga de trabajo percibida, UX, hardware complejo, \textit{digital model}, \textit{digital shadow}, \textit{digital twin} e interoperabilidad; selección del caso de estudio; definición de metas de rendimiento y alcance funcional.

\paragraph{Fase II: Pipeline de Activos y Visualización} Limpieza, optimización e integración del modelo CAD; preparación de materiales y mallas para visualización en tiempo real; configuración del proyecto Unity y del \textit{pipeline} web.

\paragraph{Fase III: Implementación del Prototipo} Desarrollo de interacciones, selección de piezas, flujo de detalle, modos de análisis visual, vista explosionada, corte transversal e interfaces responsivas.

\paragraph{Fase IV: Evaluación Formativa y Cierre} Medición de KPIs técnicos, aplicación de SUS, NASA-TLX Raw y \textit{Think-Aloud}, consolidación del informe, manuales y anexos, y formulación de un roadmap de evolución hacia \textit{twin manifest}, telemetría histórica, \textit{live digital shadow} y modo de servicio (véase la Figura \ref{fig:fases_trabajo_proyecto}).

\begin{figure}[htbp]
    \caption{Fases de Trabajo del Proyecto}
    \label{fig:fases_trabajo_proyecto}
    \includegraphics[width=0.94\textwidth,height=0.45\textheight,keepaspectratio]{figures/chapter3/fig_3_fases_trabajo.pdf}
    \apanote[11.5pt]{Las fases agrupan fundamentación, preparación de activos, implementación y cierre evaluativo-documental. Elaboración propia.}
\end{figure}

\subsection{Prueba Breve Formativa Previa a la Evaluación Final}

En este informe, las etiquetas P1 y P2 se usan como niveles internos de prioridad de QA: P1 agrupa correcciones críticas para que el flujo principal sea usable y defendible; P2 agrupa ajustes importantes de pulido, consistencia o reducción de fricción que no bloquean por sí solos la demostración del prototipo.

Antes de aplicar los instrumentos con la muestra final, se realizó una prueba breve exploratoria con usuarios para observar fricciones tempranas de interacción. Esta actividad se considera una iteración formativa previa, alineada con el enfoque \textit{Think-Aloud} y con observación de acciones espontáneas o inconscientes, y permitió ajustar la versión que luego fue evaluada con SUS, NASA-TLX Raw y protocolo comparativo 3D/2D.

El objetivo de esta prueba previa fue identificar gestos que los usuarios intentaban realizar de manera automática antes de recibir instrucción detallada. Entre los patrones observados se registraron intentos de hacer \textit{swipe} sobre las tarjetas del tutorial, uso de gestos verticales para abrir o cerrar el panel de información, sensibilidad excesiva percibida en el zoom táctil, expectativa de doble toque para aislar elementos y rechazo a la apertura automática del panel de información cuando la intención principal era aislar una pieza. Estos hallazgos alimentaron los ajustes de cierre de la interfaz para hacer que la aplicación respondiera mejor a expectativas naturales de uso en la evaluación posterior.

La revisión formativa produjo una lista de correcciones priorizadas que se trabajó antes del cierre de QA de P1 y P2. En particular, se ajustaron los gestos del tutorial, la apertura y cierre del \textit{bottom sheet}, el bloqueo táctil del panel de información para evitar interacción accidental con el dron, la sensibilidad de zoom y órbita, el doble toque sobre piezas y hotspots, la navegación por capas con el botón Atrás del navegador, la eliminación de la X externa del contenedor WebGL y la estabilización de iconos y microinteracciones en móvil. Por tanto, esta prueba breve se reporta como evidencia formativa de diseño iterativo y se distingue de los resultados con usuarios consolidados en el capítulo 5 (véase la Figura \ref{fig:ciclo_qa_formativa}).

\begin{figure}[htbp]
    \caption{Ciclo de QA Formativo y Optimización de Interacción}
    \label{fig:ciclo_qa_formativa}
    \includegraphics[width=0.88\textwidth,height=0.45\textheight,keepaspectratio]{figures/chapter3/fig_3_ciclo_qa_formativa.pdf}
    \apanote[-18.3pt]{Proceso iterativo de detección de fricciones de interfaz y su correspondiente ajuste funcional en la build. Elaboración propia.}
\end{figure}

\FloatBarrier
\subsection{Muestra y Perfil de Participantes}

Se proyectó una muestra no probabilística por conveniencia de al menos 30 participantes como escenario deseable para fortalecer la lectura descriptiva de los instrumentos estandarizados.

El estudio integrado en el informe se realizó con 12 participantes anonimizados, dentro del escenario mínimo operativo definido para evaluación formativa. El perfil esperado corresponde a estudiantes avanzados, técnicos o profesionales de ingeniería, manufactura digital o mantenimiento. Esta decisión se apoyó en la tradición de evaluaciones formativas en HCI y UX, donde muestras reducidas permiten identificar problemas recurrentes de interacción y generar insumos de rediseño (Nielsen, 2000; Lazar et al., 2017). En consecuencia, los resultados de SUS y NASA-TLX Raw se interpretan de forma descriptiva y exploratoria, sin pretensión de inferencia estadística poblacional.

La caracterización efectiva de la muestra se reporta de forma agregada para preservar el anonimato de los participantes y evitar la asociación directa entre perfil profesional y respuestas individuales. La muestra quedó conformada por perfiles afines a ingeniería, desarrollo de software, diseño de interacción y visualización 3D, coherentes con el carácter técnico-formativo de la evaluación (véase la Tabla \ref{tab:perfil_agregado_participantes}).

\begin{table}[htbp]
    \caption{Perfil Académico o Profesional Agregado de la Muestra}
    \label{tab:perfil_agregado_participantes}
    \begin{tabularx}{\textwidth}{Xr}
        \hline
        \textit{Perfil académico o profesional} & \textit{Cantidad} \\ \hline
        Ingeniería de sistemas, estudiantes o profesionales & 3 \\
        Desarrollo de software & 1 \\
        Ingeniería multimedia & 1 \\
        Ingeniería mecatrónica & 1 \\
        Ingeniería civil & 1 \\
        Ingeniería electrónica & 2 \\
        Ingeniería mecánica & 1 \\
        Diseño gráfico con énfasis en UX/UI & 1 \\
        Entusiasta 3D & 1 \\ \hline
        \textit{Total} & \textit{12} \\ \hline
    \end{tabularx}
    \apanote[-18.3pt]{La tabla presenta datos agregados de caracterización y no permite asociar perfiles con respuestas individuales. Elaboración propia.}
\end{table}

\subsection{Variables de Investigación}

\subsubsection{Variable Independiente}

\paragraph{Tipo de Visualización Técnica} Medio utilizado para consultar y comprender la información del sistema. Su definición operacional compara el visor 3D interactivo en web frente a un soporte 2D de referencia en dos niveles (visor 3D interactivo / documentación 2D).

\subsubsection{Variables Dependientes}

\paragraph{Usabilidad Percibida} Grado en que el usuario percibe el sistema como fácil de aprender, consistente y útil, medido mediante el puntaje total del cuestionario \textit{System Usability Scale} (Brooke, 1996) en escala de 0 a 100. El puntaje 68 se toma como referencia del promedio histórico del instrumento (Bangor et al., 2008, 2009; Sauro \& Lewis, 2016).

\paragraph{Carga de Trabajo Percibida} Percepción subjetiva de demanda y esfuerzo durante la ejecución de tareas, correspondiente al promedio de las seis dimensiones del NASA-TLX Raw (Hart \& Staveland, 1988), en escala de 0 a 100.

\paragraph{Desempeño en Tareas de Inspección} Capacidad del usuario para completar tareas representativas de identificación, ubicación y consulta técnica sobre el ensamblaje, medida en tiempo (segundos para T1--T3), tasa de tareas completadas, completitud autónoma, errores observables y ayudas requeridas.

\paragraph{Rendimiento Técnico y Uso de Memoria} Desempeño del prototipo web durante la carga y la interacción en tiempo real (FPS, \textit{frame time}, \textit{draw calls}, presupuesto geométrico, tamaño de build, \textit{heap} de WebAssembly) bajo metas operativas de 30 FPS o más y \textit{frame time} $\leq 33{,}33$\,ms en entorno controlado.

\subsubsection{Variables de Control y Condiciones Registradas}

Las condiciones de control registradas incluyen: código anónimo del participante, fecha, perfil profesional agregado, experiencia previa, dispositivo, sistema operativo, navegador y resolución; condición aplicada (visor 3D o soporte 2D); orden contrabalanceado de exposición (AB / BA); y versión de build, estado de caché y archivos exportados por el profiler interno.

El entorno de prueba controlado quedó delimitado por la misma build usada en la recolección de datos técnicos y en la evaluación con usuarios. Para la copia pública de cierre, la trazabilidad se fija en la rama \texttt{master} del repositorio y en los artefactos versionados en \texttt{docs/Build}, \texttt{Informe\_final/} y \texttt{Telemetria/Mediciones\_WebGL}. La build se ejecutó con la caché del navegador previamente cargada. En escritorio se usó el mismo equipo empleado en las pruebas técnicas: Intel Core i7-5820K, NVIDIA GTX 980 Ti, 48 GB RAM, resolución efectiva del viewport de 1910\(\times\)903, Windows 11 y Google Chrome 141. En móvil se utilizó un Redmi Note 10S con MIUI Global 14.0.11 y Google Chrome 148. Estas condiciones se registran para asegurar trazabilidad y evitar extrapolar el comportamiento observado a equipos o configuraciones no evaluadas. Cuando los archivos exportados por Unity WebGL reportan sistemas como \texttt{Windows 10} o \texttt{Android 10}, ese valor se interpreta como metadato normalizado del navegador o del runtime, y se prioriza el registro manual del dispositivo, sistema operativo y navegador consignado por el moderador.

\subsection{Diseño Comparativo y Condiciones de Prueba}

La evaluación con usuarios se planteó como un diseño intra-sujeto de carácter formativo. Cada participante interactuó con dos condiciones funcionalmente equivalentes: \textbf{Condición A}, correspondiente al prototipo 3D interactivo en navegador, y \textbf{Condición B}, correspondiente a un soporte 2D de referencia construido a partir de documentación técnica estática del mismo sistema. La comparación no pretendió probar causalidad psicológica concluyente, sino describir diferencias observables de percepción y uso entre ambos medios.

\paragraph{Definición Operativa de la Condición 2D} La condición 2D se implementó mediante un paquete documental controlado, fijo para todos los participantes, derivado de los manuales del fabricante ---\textit{X500 Kit Assembly Manual} y \textit{Holybro X500 Frame Kit Assembly Guide}--- y complementado con una lámina de rotulado alineada con la taxonomía canónica del proyecto (Holybro, s.\,f.-a, s.\,f.-b).

Para mantener equivalencia informativa con la condición 3D, dicho paquete incluyó una vista general del ensamblaje, una secuencia o despiece del frame, una lámina de identificación de subsistemas de propulsión y electrónica, y una ficha textual de componentes por categoría. Durante la prueba, los participantes consultaron únicamente ese paquete estático mediante lectura, desplazamiento o cambio de página, sin rotación 3D, filtros interactivos ni herramientas analíticas propias de la condición 3D.

Para mitigar sesgos de aprendizaje o fatiga, el orden de exposición se contrabalanceó mediante dos secuencias posibles: AB y BA. La asignación del orden se realizó por alternancia simple: los participantes con código impar siguieron la secuencia AB, es decir, visor 3D primero y soporte 2D después; los participantes con código par siguieron la secuencia BA, es decir, soporte 2D primero y visor 3D después. Esta regla permitió distribuir de forma controlada el posible efecto de aprendizaje o fatiga sin introducir una asignación subjetiva por parte del moderador. Las tareas mantuvieron el mismo objetivo funcional en ambas condiciones y variaron únicamente en el medio de representación. El moderador empleó instrucciones equivalentes y una misma matriz de registro para tiempo en T1--T3, éxito, errores, ayudas y observaciones.

El entorno técnico se documentó de forma explícita en las sesiones de rendimiento WebGL, incluyendo equipo, sistema operativo, navegador, resolución, escenario evaluado y archivo exportado por el profiler interno. Los KPIs técnicos de la build web se interpretan junto con esas condiciones de ejecución para evitar extrapolaciones ambiguas del desempeño.

En cuanto al alcance de plataforma, el prototipo se concibió para ejecutarse en navegadores web compatibles de escritorio y móviles soportados por la plataforma web de Unity. La compatibilidad efectiva depende del navegador, su versión, la memoria disponible y el nivel de optimización de la build. Por ello, la evaluación técnica registró dispositivo, navegador, resolución, escenario y medidas fechadas, evitando extrapolar el comportamiento observado a equipos no incluidos en la prueba (Unity Technologies, s.\,f.).

Una mejora operativa de la condición 3D frente a la referencia 2D se define por la convergencia descriptiva de varios indicadores: mayor porcentaje de éxito en tareas, menos errores o ayudas requeridas, carga de trabajo percibida igual o inferior y verbalizaciones cualitativas que indiquen mejor orientación durante la inspección. El SUS se utiliza por separado como lectura global de usabilidad del prototipo 3D y no como métrica comparativa entre condiciones (véase la Figura \ref{fig:diseno_comparativo_3d_2d}).

\begin{figure}[htbp]
    \caption[Diseño Comparativo 3D Frente a 2D]{Diseño Comparativo entre Condición 3D Interactiva y Soporte 2D de Referencia}
    \label{fig:diseno_comparativo_3d_2d}
    \includegraphics[width=\textwidth,height=0.45\textheight,keepaspectratio]{figures/chapter3/fig_3_diseno_comparativo_3d_2d.pdf}
    \apanote[-18.3pt]{El diseño usa tareas equivalentes y orden contrabalanceado; SUS se reporta para el prototipo 3D como lectura global de usabilidad. Elaboración propia.}
\end{figure}

\subsection{Instrumentos}
\label{sec:instrumentos}

\paragraph{KPIs Técnicos} Unity Profiler mediante \textit{Development Build} y \textit{Autoconnect Profiler} cuando corresponda, herramientas del navegador, auditorías de cobertura y jerarquía, y el profiler interno de la aplicación WebGL para capturar sesiones por escenario desde la propia build.

\paragraph{Registro de Desempeño en Tareas} Matriz estructurada por condición para documentar tiempo en segundos para T1, T2 y T3, porcentaje de éxito, completitud autónoma, errores observables, ayudas requeridas y observaciones.

\paragraph{SUS} Cuestionario posterior a la interacción con el prototipo 3D para estimar usabilidad percibida mediante 10 ítems en escala Likert de 1 a 5.

\paragraph{NASA-TLX Raw} Cuestionario posterior a tareas para estimar carga de trabajo percibida en seis dimensiones, cada una en escala de 0 a 100. Se usa como medida de carga de trabajo percibida y no como medición directa de la Teoría de la Carga Cognitiva. En este informe se adopta la variante \textit{Raw TLX} por su pertinencia para lectura formativa rápida sin ponderación pareada adicional (Hart \& Staveland, 1988; Hart, 2006).

\paragraph{Think-Aloud} Protocolo de verbalización concurrente con instrucción estandarizada, intervención neutral del moderador y apoyo de matriz de observación, audio y captura de pantalla cuando el entorno lo permita (véase la Figura \ref{fig:instrumentos_evaluacion_prototipo}).

\begin{figure}[htbp]
    \caption{Instrumentos de Evaluación del Prototipo}
    \label{fig:instrumentos_evaluacion_prototipo}
    \includegraphics[width=0.92\textwidth,height=0.45\textheight,keepaspectratio]{figures/chapter3/fig_3_instrumentos_evaluacion.pdf}
    \apanote[+6.2pt]{Los instrumentos se organizan por capas técnica, conductual, perceptiva y cualitativa para evitar depender de un único indicador. Elaboración propia.}
\end{figure}

\FloatBarrier
\subsection{Protocolo de Tareas}

Las tareas de validación se alinearon con el flujo visible final de la aplicación y se organizaron en cuatro bloques equivalentes para ambas condiciones:

\paragraph{Bloque 1 (Ubicación de Componente)} Localizar un componente específico del dron en el visor 3D o en el soporte 2D.

\paragraph{Bloque 2 (Interpretación Funcional)} Explicar la función general del componente y su relación con el sistema.

\paragraph{Bloque 3 (Relación Espacial)} Identificar vínculos espaciales entre el componente, piezas cercanas y el ensamblaje.

\paragraph{Bloque 4 (Análisis Visual Guiado)} Usar una herramienta visible del visor ---por ejemplo \texttt{Explode}, \texttt{Cut}, \textit{X-Ray}, \textit{Thermal} o presets de \textit{Studio}--- o interpretar la información equivalente en el soporte 2D.

\subsection{Procedimiento de Análisis}

\subsubsection{Análisis Cuantitativo}

El componente cuantitativo del estudio tuvo carácter descriptivo y exploratorio. La lectura cuantitativa se organizó por tipo de dato: desempeño en tareas por condición, usabilidad global del prototipo 3D, carga de trabajo percibida por condición y KPIs técnicos de la build web bajo entorno controlado.

El puntaje SUS se calcula según el procedimiento estándar de Brooke:
\[
SUS = \left(\sum_{j \in impares}(R_j - 1) + \sum_{j \in pares}(5 - R_j)\right)\times 2.5
\]
donde $R_j$ corresponde a la respuesta del participante en el ítem $j$. Un puntaje alto indica mejor usabilidad percibida. Para su lectura, 68 se considera una referencia histórica aproximada del instrumento, mientras que puntajes en el rango 70--72 se interpretan como favorables o por encima del centro de referencia y los valores en los ochenta como claramente superiores, siempre de forma descriptiva y no como umbral absoluto (Bangor et al., 2008, 2009; Sauro \& Lewis, 2016). En este estudio, el SUS se reporta únicamente para la condición 3D.

Para NASA-TLX Raw se promedian las seis dimensiones del instrumento y se obtiene un resultado independiente para cada condición de prueba:
\[
NASA\text{-}TLX_{Raw} = \frac{1}{6}\sum_{i=1}^{6} D_i
\]
donde $D_i$ corresponde al puntaje observado en demanda mental, demanda física, demanda temporal, rendimiento, esfuerzo y frustración. En esta adaptación se utiliza la variante \textit{Raw TLX}, por lo que no se aplica ponderación por comparaciones pareadas. La dimensión de rendimiento se diligencia con orientación invertida (0 = desempeño perfecto, 100 = desempeño deficiente) para conservar coherencia direccional en el promedio global. Además del promedio agregado, se reportan las subescalas para evitar interpretaciones excesivamente reductivas (Hart \& Staveland, 1988; Hart, 2006).

Los KPIs técnicos se contrastaron con las metas operativas del proyecto a partir de datos obtenidos sobre la build publicada. La aplicación integra herramientas internas de captura por escenario en JSON/CSV; por ello el informe reporta FPS, \textit{frame time}, memoria, artefactos de build, presupuesto geométrico y compatibilidad mediante tablas cuantitativas trazables.

\subsubsection{Análisis Cualitativo}

El protocolo \textit{Think-Aloud} se aplicó en modalidad concurrente durante la ejecución de tareas. A cada participante se le solicitó verbalizar qué intentaba hacer, cómo interpretaba la interfaz, qué dudas le surgían, qué errores percibía y con qué criterio decidía su siguiente acción. El evaluador registró verbalizaciones relevantes, pausas, puntos de vacilación, errores de navegación, solicitudes de ayuda y estrategias de recuperación mediante una matriz estructurada por tarea y por condición.

Posteriormente, las observaciones se codificaron por categorías analíticas: comprensión espacial del ensamblaje, navegación y control de cámara, interpretación de la interfaz, correspondencia terminológica, errores o bloqueos y sugerencias de mejora. La codificación se realizó mediante lectura de los registros del moderador, una segunda pasada de revisión y consolidación de memos breves por participante. Esta lectura cualitativa permite triangular los resultados de SUS y NASA-TLX y explicar por qué determinadas tareas resultaron más claras, ambiguas o demandantes para el usuario (Ericsson \& Simon, 1993).

\subsubsection{Triangulación de Hallazgos}

La integración final de resultados sigue una lógica de triangulación convergente. En primer lugar, los indicadores de desempeño en tareas y los KPIs técnicos describen la viabilidad operativa de la build web en el entorno de prueba documentado. En segundo lugar, SUS aporta una lectura global de usabilidad del prototipo 3D, mientras que NASA-TLX compara la carga de trabajo percibida entre la condición interactiva y la referencia 2D. En tercer lugar, el \textit{Think-Aloud} y la observación estructurada explican por qué determinadas tareas producen fricción, claridad o ambigüedad. Un hallazgo se considera más sólido cuando muestra consistencia entre estas capas de evidencia, teniendo en cuenta que la comparación entre 3D y 2D se apoya en desempeño y carga de trabajo percibida, mientras que la usabilidad global se informa para el prototipo interactivo (véase la Figura \ref{fig:triangulacion_evidencia}).

\begin{figure}[htbp]
    \caption{Triangulación de Evidencia para Interpretar el Prototipo}
    \label{fig:triangulacion_evidencia}
    \includegraphics[width=0.92\textwidth,height=0.45\textheight,keepaspectratio]{figures/chapter3/fig_3_triangulacion_evidencia.pdf}
    \apanote[+6.2pt]{La interpretación final combina KPIs técnicos, desempeño en tareas, SUS, NASA-TLX Raw, \textit{Think-Aloud} y observación estructurada. Elaboración propia.}
\end{figure}

\FloatBarrier
\subsection{Consideraciones Éticas}

La participación en pruebas fue voluntaria. Los participantes fueron identificados mediante códigos anónimos y los resultados se reportan de forma agregada o anonimizada. No se recolectaron datos sensibles ni se realizaron intervenciones físicas. Cualquier registro de audio, pantalla o fotografía quedó condicionado a autorización explícita, y los documentos de validación se conservaron únicamente como soporte académico del proyecto. El estudio no contempla exposición a riesgos adicionales para los participantes.

\input{chapters/04_desarrollo}
\input{chapters/05_resultados}
% =============================================================================
% Capítulo 6 — Conclusiones y Recomendaciones
% Conforme a APA 7 UNAD (Instructivo marzo 2023)
% =============================================================================

\section{Conclusiones}

El proyecto permitió construir y validar de forma descriptiva TwinSight X500 como prototipo WebGL para visualización técnica de hardware complejo. La build publicada integra el flujo principal de exploración, selección, aislamiento, ficha contextual, modos de inspección/análisis/estudio, corte transversal, vista explosionada, filtros y visualización térmica heurística. Los resultados técnicos y de usuario se integraron con datos fechados: mediciones del profiler interno, reducción geométrica del activo y evaluación con 12 participantes mediante SUS, NASA-TLX Raw y protocolo \textit{Think-Aloud}. Por su alcance real, el artefacto queda delimitado como \textit{visual product twin}: un modelo visual-semántico de producto preparado para inspección técnica y evolución futura, sin telemetría en vivo ni sincronización operacional con un activo físico.

\subsection{Respecto al Objetivo General}

Se logró construir un visor web interactivo para la visualización técnica del Holybro X500~V2, accesible desde navegador y capaz de integrar navegación 3D, selección de piezas, información contextual y herramientas de análisis visual. La validación formativa con 12 participantes registró un puntaje SUS promedio de 91,88, un promedio NASA-TLX Raw de 8,69 para el visor 3D frente a 19,89 para el soporte 2D y menor tiempo medio en T1, T2 y T3 para la condición 3D. Estos resultados se interpretan de forma descriptiva, dentro de una muestra no probabilística y un diseño formativo, y muestran una recepción favorable del prototipo junto con una menor carga de trabajo percibida durante la exploración 3D. Además, la taxonomía por piezas, los grupos de escena, las fichas y las herramientas de inspección constituyen una base razonable para evolucionar hacia modelos con datos operativos, siempre que esa evolución se aborde por fases.

\subsection{Respecto a los Objetivos Específicos}

\paragraph{OE1: Pipeline de Optimización y Preparación de Activos} Se consolidaron el caso de estudio, la taxonomía canónica y un pipeline defendible que combina importación CAD, saneamiento geométrico, remodelado hard-surface, modularización delimitada de fasteners, automatización puntual y preparación para runtime web. La comparación geométrica muestra una reducción desde rutas CAD superiores a 6,5 millones de triángulos hasta un modelo base optimizado de 95\,617 triángulos, incluyendo piezas adicionales modeladas desde cero como hélices. Esta reducción permitió conservar legibilidad formal y hacer viable la interacción WebGL dentro del alcance del MVP, sin confundir ese conteo de activo base con la complejidad de la escena runtime instrumentada.

\paragraph{OE2: Materiales, Modos de Visualización y Rendimiento WebGL} Se integró en Unity URP una capa visual que combina renderizado realista, modos analíticos y presets de presentación: \textit{Realistic}, \textit{X-Ray}, \textit{Solid View}, \textit{Thermal}, \textit{Blueprint}, entornos de \textit{Studio}, corte transversal, filtros y vista explosionada. Las mediciones muestran cumplimiento pleno de la meta de 30 FPS en escritorio y equipos de mayor capacidad; en móviles de gama media-alta el comportamiento fue cercano a la meta y en el límite inferior quedó por debajo de 30 FPS, aunque mantuvo navegabilidad en el flujo principal. Por ello, la compatibilidad móvil debe leerse como funcional en los dispositivos evaluados y no como garantía universal para cualquier hardware.

\paragraph{OE3: Implementación del Prototipo Interactivo} Se obtuvo una arquitectura funcional basada en Unity Web, C\#, UI Toolkit, gestores de escena, orquestación de modos y saneamiento del import mediante \texttt{ImportedDroneRuntimeBinder}. La build publicada integra navegación, selección, aislamiento, ficha contextual inferior, menús \textit{Inspect}, \textit{Analyze} y \textit{Studio}, \texttt{Explode}, \texttt{Cut}, filtros, modos visuales, tutorial y adaptación de la experiencia \textit{mobile-first} a escritorio. La separación entre UI visible, código oculto y código no integrado mejora la trazabilidad del sistema y evita documentar como alcance final capacidades que permanecen fuera del flujo publicado.

\paragraph{OE4: Evaluación Formativa 3D vs.\ 2D, Usabilidad y Carga de Trabajo} La evaluación con 12 participantes permitió describir la experiencia del prototipo desde capas complementarias. La matriz de desempeño registró cuatro tareas equivalentes por condición y 96 registros tarea-condición completados, con una ayuda baja en el visor 3D y sin errores de tarea en la matriz estructurada. En las tres tareas cronometradas, el visor 3D presentó menor tiempo medio que el soporte 2D, mientras que NASA-TLX Raw mostró menor carga percibida en la condición 3D en los 12 casos observados. La completitud total en ambas condiciones sugiere un efecto techo en la matriz de desempeño, por lo que la ventaja del visor 3D se interpreta principalmente en términos de menor demanda subjetiva, mejor orientación espacial y menor tiempo descriptivo de resolución, no como prueba inferencial de superioridad causal. El \textit{Think-Aloud} explicó esta lectura: los participantes valoraron la comprensión espacial, el aislamiento, las vistas analíticas y el panel de pieza, mientras que las principales fricciones se concentraron en navegación móvil, lectura de iconos y selección de elementos pequeños. La documentación producida queda como soporte de transferencia y mantenimiento, no como objetivo específico independiente.

\subsection{Limitaciones}

La evaluación con usuarios estuvo delimitada por una muestra no probabilística de 12 participantes, por lo que los resultados cuantitativos se interpretan de manera descriptiva y formativa sin pretensión de inferencia estadística poblacional. Las pruebas se ejecutaron en dispositivos controlados con la caché del navegador precargada sobre la build trazada en la rama \texttt{master}; asimismo, la T4 se diseñó como tarea exploratoria guiada y no cronometrada, concentrando el análisis temporal en T1--T3 como indicadores de desempeño relativo entre condiciones.

En el frente técnico, las mediciones de rendimiento reflejan el comportamiento observado en los dispositivos, navegadores y resoluciones registrados, constituyendo evidencia empírica situada y no una garantía universal para cualquier hardware móvil. Por restricciones de tiempo en modelado y optimización dentro de Unity, el alcance del modelo no incluyó la totalidad de cables, arneses ni componentes electrónicos microscópicos. De igual modo, la versión de escritorio opera como una adaptación funcional de la interfaz \textit{mobile-first}, y la visualización térmica responde a una aproximación heurística por componentes antes que a una simulación de elementos finitos (FEA) o termografía calibrada.

Finalmente, el artefacto se delimita conceptualmente como un \textit{visual product twin}: un modelo visual-semántico que organiza la comprensión del ensamblaje pero carece de telemetría en vivo, sincronización operacional con un activo físico, conectores IoT/SCADA/PLM o un esquema de \textit{twin manifest} interoperable. Determinadas herramientas técnicas permanecen implementadas en código como capacidades secundarias u ocultas, sin formar parte del flujo visible publicado para el usuario final.

\subsection{Trabajo Futuro}

El trabajo futuro se organiza como una ruta de madurez incremental (véase la Figura \ref{fig:roadmap_visual_product_twin}). La prioridad consiste en ampliar evidencia, datos e interoperabilidad por etapas, antes de aspirar a inteligencia artificial predictiva o a un gemelo digital operacional completo.

\paragraph{Fase 0: Ampliación del Visual Product Twin} Replicar la validación con una muestra mayor, idealmente cercana a 30 participantes, contrastar nuevos dispositivos y depurar la interacción móvil con énfasis en navegación, iconos y selección de piezas pequeñas.

\paragraph{Fase 1: Formalización del \textit{Twin Manifest}} Definir un esquema JSON o equivalente que conecte cada pieza con identificador estable, categoría, material, función, sensores posibles, variables esperadas, fallas comunes, repuestos, procedimientos y documentos asociados.

\paragraph{Fase 2: Telemetría Histórica y Digital Shadow de Reproducción} Importar logs o datasets CSV/JSON para reproducir estados por componente en una línea de tiempo, sin exigir todavía conexión en vivo.

\paragraph{Fase 3: Live Digital Shadow} Incorporar adaptadores WebSocket, MQTT, REST, MAVLink, ROS, OPC UA o W3C WoT según el dominio, con frecuencia de actualización definida y control de rendimiento WebGL.

\paragraph{Fase 4: Modo Servicio y Mantenimiento Guiado} Añadir síntomas, causas probables, checklists, herramientas, repuestos, evidencia, severidad y exportación de reportes por pieza.

\paragraph{Fase 5: Digital Twin Operacional} Integrar modelos calibrados, simulaciones \textit{what-if}, predicción de umbrales y decisiones human-in-the-loop solo cuando existan datos históricos confiables y validación con eventos reales.

\paragraph{Mejoras Transversales} Diseñar una UI desktop específica, modelar cables y conexiones electrónicas relevantes, evaluar con datos si una ruta nativa de LOD aporta una mejora real al pipeline, calibrar el subsistema térmico e incorporar soporte multi-idioma.

\begin{figure}[htbp]
    \caption{Roadmap Visual de Evolución desde Visual Product Twin hasta Digital Twin Operacional}
    \label{fig:roadmap_visual_product_twin}
    \includegraphics[width=\textwidth,height=0.45\textheight,keepaspectratio]{figures/chapter6/fig_6_roadmap_visual_product_twin.pdf}
    \apanote[14.4pt]{La ruta separa el cierre visual documentado de fases posteriores con manifest, telemetría, sombra digital en vivo, modo servicio y gemelo operacional. Elaboración propia.}
\end{figure}

\section{Recomendaciones}

Se recomienda a futuros desarrolladores e investigadores mantener una taxonomía documental única y formalmente reconciliada entre piezas canónicas, \textit{anchors} de escena y \textit{renderers} de runtime. Esta coherencia evita discrepancias entre la capa de datos semánticos y la jerarquía física de renderizado. Asimismo, se aconseja documentar como alcance final exclusivamente las funciones visibles y verificables de la build publicada, preservando registros fechados de versión, dispositivo, navegador y resolución cada vez que se ejecuten pruebas de rendimiento o sesiones con usuarios.

Por otra parte, se sugiere conservar el código modular no expuesto como soporte técnico auditable, manteniendo una distinción explícita entre componentes integrados, capacidades ocultas y desarrollos proyectados. Finalmente, se recomienda sostener la denominación rigurosa de \textit{visual product twin} para el alcance actual del visor, reservando los términos \textit{digital shadow} o \textit{digital twin} para aquellas fases que integren telemetría en tiempo real, esquemas formales de interoperabilidad, modelos físicos calibrados y validación en entornos de operación industrial.

\input{chapters/07_referencias}
% =============================================================================
% Capítulo 8 — Apéndices
% Conforme a APA 7 UNAD (Instructivo marzo 2023)
% =============================================================================

\section{Apéndices}

Este capítulo lista los artefactos anexos activos que respaldan el informe, la app publicada, la validación técnica y la evaluación con usuarios.

\newpage
\phantomsection
\addcontentsline{toc}{subsection}{Apéndice A: Repositorio de Código Fuente}
\label{apendice:A}
{\centering\textbf{Apéndice A}\par
\textit{Repositorio de Código Fuente}\par}

Repositorio principal del proyecto:
\begin{center}
\url{https://github.com/delarge95/WebGL-Thesis-Proposal}
\end{center}

Este repositorio, junto con el \textit{snapshot} local del proyecto y los anexos versionados, funciona como referencia técnica principal para escena, scripts, herramientas de editor y documentación de desarrollo. La copia pública de cierre se consolida en la rama \texttt{master}; por tanto, la trazabilidad pública debe verificarse contra el \texttt{HEAD} vigente de esa rama, la URL de demo publicada y los artefactos activos listados en este capítulo.

Si el repositorio público no expone todos los artefactos internos por razones de limpieza, peso o información sensible, la evaluación académica debe apoyarse en la URL de demo publicada, los anexos activos del informe, la documentación técnica interna versionada y el \textit{snapshot} local disponible para revisión.

\newpage
\phantomsection
\addcontentsline{toc}{subsection}{Apéndice B: Manual de Usuario}
\label{apendice:B}
{\centering\textbf{Apéndice B}\par
\textit{Manual de Usuario}\par}

Ruta del manual activo:
\begin{center}
\texttt{Informe\_final/Manual\_de\_usuario/manual\_usuario.tex}
\end{center}

\newpage
\phantomsection
\addcontentsline{toc}{subsection}{Apéndice C: Manual Técnico}
\label{apendice:C}
{\centering\textbf{Apéndice C}\par
\textit{Manual Técnico}\par}

Ruta del manual activo:
\begin{center}
\texttt{Informe\_final/Manual\_tecnico/manual\_tecnico.tex}
\end{center}

\newpage
\phantomsection
\addcontentsline{toc}{subsection}{Apéndice D: Documentación Técnica Activa}
\label{apendice:D}
{\centering\textbf{Apéndice D}\par
\textit{Documentación Técnica Activa}\par}

Documentos activos:

\path{Informe_final/Desarrollo_App/Documentacion_Tecnica/01_Arquitectura_del_Sistema.md}

\path{Informe_final/Desarrollo_App/Documentacion_Tecnica/02_Referencia_Tecnica_Modulos.md}

\path{Informe_final/Desarrollo_App/Documentacion_Tecnica/03_Pipeline_Renderizado_Shaders.md}

\path{Informe_final/Desarrollo_App/Documentacion_Tecnica/04_Guia_Despliegue.md}

\path{Informe_final/Desarrollo_App/Documentacion_Tecnica/05_Configuracion_WebGL.md}

\path{Informe_final/Desarrollo_App/Documentacion_Tecnica/08_Sistema_Termico_Hibrido.md}

La versión publicada para consulta desde la landing se conserva en \path{docs/manuals/}. Esta copia se sincroniza con los documentos internos cuando se actualiza el cierre técnico.

\newpage
\phantomsection
\addcontentsline{toc}{subsection}{Apéndice E: Matriz de Desconexiones}
\label{apendice:E}
{\centering\textbf{Apéndice E}\par
\textit{Matriz de Desconexiones}\par}

La matriz maestra de desconexiones entre app y documentación se encuentra en:
\begin{center}
\path{Informe_final/Desarrollo_App/MATRIZ_DESCONEXIONES_APP_DOCUMENTACION.md}
\end{center}

\newpage
\phantomsection
\addcontentsline{toc}{subsection}{Apéndice F: Catálogo Canónico del Holybro X500 V2}
\label{apendice:F}
{\centering\textbf{Apéndice F}\par
\textit{Catálogo Canónico del Holybro X500 V2}\par}

La taxonomía semántica del proyecto se fija a partir de:
\begin{center}
\path{desarrollo/unity_project/Assets/Resources/x500v2_info_panel_bilingual_catalog.json}
\end{center}

Esta fuente pública integra los identificadores canónicos, las fichas bilingües de información y las claves de trazabilidad utilizadas por el panel de la app. La taxonomía académica conserva 28 piezas canónicas; la escena final opera adicionalmente con 30 \textit{anchors} y 257 \textit{renderers}/\textit{colliders} auditados.

\newpage
\phantomsection
\addcontentsline{toc}{subsection}{Apéndice G: Validación Funcional de la Escena Final}
\label{apendice:G}
{\centering\textbf{Apéndice G}\par
\textit{Validación Funcional de la Escena Final}\par}

El documento de referencia para la build funcional es:
\begin{center}
\path{Informe_final/Desarrollo_App/VALIDACION_FUNCIONAL_FINA_2026-04-09.md}
\end{center}

\newpage
\phantomsection
\addcontentsline{toc}{subsection}{Apéndice H: Instrumentos y Registros de Evaluación}
\label{apendice:H}
{\centering\textbf{Apéndice H}\par
\textit{Instrumentos y Registros de Evaluación}\par}

\setcounter{table}{0}
\renewcommand{\thetable}{H\arabic{table}}
\setcounter{figure}{0}
\renewcommand{\thefigure}{H\arabic{figure}}

Los instrumentos metodológicos activos del cierre son el cuestionario SUS, el NASA-TLX Raw, el protocolo \textit{Think-Aloud}, la guía de mediciones técnicas WebGL, la guía de tareas comparativas 3D vs.\ 2D y las tablas de KPIs técnicos asociadas a sesiones fechadas del profiler interno.

Las tareas de prueba se alinean con cuatro acciones equivalentes: ubicar un componente, interpretar su función, explicar una relación espacial y resolver una acción de análisis visual guiado. En la condición 3D estas tareas se apoyan en navegación, selección, \textit{bottom sheet}, \textit{Explode}, \texttt{Cut} y modos visuales; en la condición 2D se resuelven mediante el paquete documental estático de referencia.

Para mantener coherencia con la metodología vigente, SUS se reporta como lectura global del prototipo 3D y NASA-TLX Raw se reporta como comparación de carga de trabajo percibida entre la condición 3D y el soporte 2D de referencia.

\begin{table}[htbp]
    \caption{Matriz de Equivalencia Informativa entre Visor 3D y Soporte 2D}
    \label{tab:equivalencia_soporte_2d_visor_3d}
    \begin{tabular}{p{3.2cm}p{4.4cm}p{4.4cm}}
        \hline
        \textbf{Tarea} & \textbf{Condición 3D} & \textbf{Condición 2D} \\ \hline
        Ubicación & Navegación orbit/pan/zoom, selección y foco. & Vista general, lámina o página del soporte. \\
        Interpretación & Ficha inferior de pieza y categoría. & Ficha textual o fragmento documental 2D. \\
        Relación espacial & Inspección 3D, aislamiento, explode o corte. & Despiece o diagrama estático. \\
        Análisis guiado & Uso de explode, cut, X-Ray, thermal o studio. & Interpretación de vista estática equivalente. \\ \hline
    \end{tabular}
    \apanote{La equivalencia es funcional, no formal: el visor 3D ofrece interacción espacial, mientras que el soporte 2D conserva información suficiente para resolver la misma intención de tarea. Elaboración propia.}
\end{table}

Las versiones activas y versionadas de instrumentos y registros se conservan en:

\path{Informe_final/validacion/CUESTIONARIO_SUS.md}

\path{Informe_final/validacion/CUESTIONARIO_NASA_TLX.md}

\path{Informe_final/validacion/GUIA_TAREAS_VALIDACION.md}

\path{Informe_final/validacion/PROTOCOLO_THINK_ALOUD.md}

\path{Informe_final/validacion/06_GUIA_MEDICIONES_TECNICAS_WEBGL.md}

\path{Informe_final/validacion/07_TABLAS_RENDIMIENTO_WEBGL_MEDICIONES.tex}

Los registros completos por participante se conservan como evidencia interna de evaluación y no forman parte de la superficie pública del repositorio. En el capítulo 5 se reportan únicamente datos agregados y participantes anonimizados; si un jurado o asesor requiere auditar los registros individuales, estos deben entregarse por un canal académico controlado y no mediante publicación pública.

Como respaldo visual del proceso de medición técnica, la Figura \ref{fig:evidencia_capturas_profiler_interno} muestra una captura puntual del profiler interno conservada junto con las capturas de cierre de la aplicación.

\begin{figure}[htbp]
    \caption{Muestra de una Captura del Profiler Interno}
    \label{fig:evidencia_capturas_profiler_interno}
    \includegraphics[width=0.92\textwidth,height=0.45\textheight,keepaspectratio]{figures/screenshots_contextual/fig_profiler_internal_evidence.png}
    \apanote{La figura documenta una captura puntual del profiler interno integrada como evidencia visual en \path{Informe_final/figures/screenshots_contextual/fig_profiler_internal_evidence.png}. Su función es respaldar visualmente el proceso de recolección técnica; los valores oficiales de FPS, memoria, renderers, mallas y triángulos reportados en las tablas técnicas provienen de las exportaciones JSON/CSV del profiler interno, no de lectura manual de esta imagen. Elaboración propia.}
\end{figure}

\newpage
\phantomsection
\addcontentsline{toc}{subsection}{Apéndice I: Estado del Pipeline de Activos}
\label{apendice:I}
{\centering\textbf{Apéndice I}\par
\textit{Estado del Pipeline de Activos}\par}

\setcounter{table}{0}
\renewcommand{\thetable}{I\arabic{table}}

Para efectos de cierre documental:

Blender y Unity conforman la ruta defendible de la versión final activa.

El FBX runtime final se documenta como unión de masters e instancias; su trazabilidad exige reporte de importación Unity, revisión de grupos, fasteners, hotspots, filtros y hélices.

Cualquier ruta nativa de LOD se documenta como decisión evaluada; si existen herramientas o scripts posteriores de LOD, no se presentan como mejora validada del estudio mientras no exista profiling y QA comparativo sobre la build publicada.

Houdini se mantiene como herramienta evaluada o experimental, no como pipeline oficial de la build cerrada, salvo evidencia nueva derivada del reimport final.

Las texturas/materiales no publicados deben mantenerse como mejora futura o documentarse expresamente si se integran a una build posterior.

El estado actual del pipeline puede verse consolidado en la Tabla \ref{tab:pipeline_3d_cierre}.

\begin{table}[htbp]
    \caption{Estado del Pipeline 3D al Cierre Documental}
    \label{tab:pipeline_3d_cierre}
    \begin{tabular}{p{2.8cm}p{2.8cm}p{2.2cm}p{4.2cm}}
        \hline
        \textbf{Etapa} & \textbf{Herramienta} & \textbf{Estado} & \textbf{Lectura editorial} \\ \hline
        Limpieza & Blender & Activa & Ruta defendible del cierre para ajuste de malla. \\
        Conversión CAD & STEPper / MoI3D & Activa / evaluada & Combinación de apertura rápida y teselación. \\
        Materiales & Unity URP & Activa & Flujo material visible en la build final. \\
        Integración & Unity & Activa & Núcleo del producto final y punto de verdad. \\
        Auditoría & Binder runtime & Activa & Coherencia entre escena y documentación. \\
        Nativa LOD & Investigación & Evaluada & Evolución de planeación técnica. \\
        Houdini & Evaluación & No canónico & Citado como exploración o trabajo futuro. \\ \hline
    \end{tabular}
    \apanote{Esta tabla evita presentar como pipeline oficial herramientas que no quedaron integradas de forma cerrada en la build final. Elaboración propia.}
\end{table}

\newpage
\phantomsection
\addcontentsline{toc}{subsection}{Apéndice J: Demo Público}
\label{apendice:J}
{\centering\textbf{Apéndice J}\par
\textit{Demo Público}\par}

La URL pública del demo publicado para la entrega es:
\begin{center}
\url{https://delarge95.github.io/WebGL-Thesis-Proposal/}
\end{center}

Esta URL corresponde al punto de acceso a la build WebGL publicada. Las métricas del capítulo 5 registran la fecha y el contexto de medición utilizado para evitar inconsistencias entre documento y producto.

\newpage
\phantomsection
\addcontentsline{toc}{subsection}{Apéndice K: Audio y Alcance No Integrado}
\label{apendice:K}
{\centering\textbf{Apéndice K}\par
\textit{Audio y Alcance No Integrado}\par}

Al momento de este cierre documental:

No existen clips de audio integrados en \texttt{Assets}.

Por defecto, el audio se clasifica como trabajo futuro si no entra en la build publicada.

El capítulo 5 reporta la evaluación técnica y de usuarios con los datos disponibles para esta versión del informe.

\newpage
\phantomsection
\addcontentsline{toc}{subsection}{Apéndice L: Auditoría Matemática y Trazabilidad con Wolfram}
\label{apendice:L}
{\centering\textbf{Apéndice L}\par
\textit{Auditoría Matemática y Trazabilidad con Wolfram}\par}

La auditoría activa de coherencia entre fórmulas documentadas, lógica de código y arquitectura real se consolida en:
\begin{center}
\path{Informe_final/Desarrollo_App/AUDITORIA_MATEMATICA_Y_ARQUITECTURA_2026-04-12.md}
\end{center}

Como soporte técnico adicional, el subsistema térmico se documenta en el anexo técnico público:
\begin{center}
\path{Informe_final/Desarrollo_App/Documentacion_Tecnica/08_Sistema_Termico_Hibrido.md}
\end{center}

Ambos documentos respaldan la interpretación correcta del modo \textit{Thermal} como simulación híbrida heurística y no como FEA en tiempo real. Los registros históricos usados durante el desarrollo permanecen como evidencia interna cuando no forman parte de la superficie pública del repositorio.

\newpage
\phantomsection
\addcontentsline{toc}{subsection}{Apéndice M: Correlación entre la Propuesta y el Informe Final}
\label{apendice:M}
{\centering\textbf{Apéndice M}\par
\textit{Correlación entre la Propuesta y el Informe Final}\par}

\setcounter{table}{0}
\renewcommand{\thetable}{M\arabic{table}}

La siguiente tabla resume la correlación de las tablas presentadas en la propuesta (véase la Tabla \ref{tab:correlacion_tablas_propuesta_informe}).

\begin{table}[htbp]
    \caption{Correlación entre Tablas de la Propuesta y el Informe Final}
    \label{tab:correlacion_tablas_propuesta_informe}
    \begin{tabular}{p{3.8cm}p{3.4cm}p{4.8cm}}
        \hline
        \textbf{Tabla en la propuesta} & \textbf{Estado en informe} & \textbf{Justificación} \\ \hline
        Comparativa Web 3D & Capítulo 2 & Sustento del estado del arte y selección de Unity. \\
        Cronograma general & Solo en propuesta & Artefacto de planeación; pierde valor en informe final. \\
        Presupuesto estimado & Solo en propuesta & Viabilidad ex ante; no es un artefacto técnico. \\
        Resultados esperados & Tabla entregados vs. esperados & Requiere contrastar expectativa inicial con entrega real. \\ \hline
    \end{tabular}
    \apanote{La ausencia de una tabla del documento final no implica incoherencia si su función era exclusivamente prospectiva en la propuesta. Elaboración propia.}
\end{table}

Asimismo, la Tabla \ref{tab:productos_entregados_esperados} presenta un contraste entre los productos entregados y los originalmente esperados.

\begin{table}[htbp]
    \caption{Productos Entregados vs.\ Esperados al Cierre Documental}
    \label{tab:productos_entregados_esperados}
    \begin{tabular}{p{3.8cm}p{2.4cm}p{5.8cm}}
        \hline
        \textbf{Producto esperado} & \textbf{Estado} & \textbf{Observación de cierre} \\ \hline
        Software web funcional & Cerrado & Flujo funcional publicado y documentable; métricas agregadas. \\
        Shaders URP optimizados & Entregado & Modos visibles/ocultos documentados de forma diferenciada. \\
        Modelos 3D optimizados & Entregado & Pipeline defendible y escena funcional para MVP. \\
        Documento de tesis & En cierre & Capítulos saneados con datos reales de rendimiento y usuarios. \\
        Informe de usabilidad & Entregado & Capítulo 5 consolida SUS, NASA-TLX Raw y Think-Aloud. \\ \hline
    \end{tabular}
    \apanote{Esta tabla actualiza el horizonte de productos de la propuesta sin borrar la diferencia entre expectativa inicial y estado real de cierre. Elaboración propia.}
\end{table}

\newpage
\phantomsection
\addcontentsline{toc}{subsection}{Apéndice N: Resumen de Artefactos Activos del Cierre}
\label{apendice:N}
{\centering\textbf{Apéndice N}\par
\textit{Resumen de Artefactos Activos del Cierre}\par}

\setcounter{table}{0}
\renewcommand{\thetable}{N\arabic{table}}

A continuación se presentan los artefactos activos resultantes del cierre (véase la Tabla \ref{tab:resumen_artefactos_cierre}).

\begin{table}[htbp]
    \caption{Resumen de Artefactos Activos del Cierre Documental}
    \label{tab:resumen_artefactos_cierre}
    \begin{tabular}{p{3.5cm}p{3.5cm}p{1.8cm}p{3.2cm}}
        \hline
        \textbf{Artefacto} & \textbf{Ubicación principal} & \textbf{Estado} & \textbf{Uso en defensa} \\ \hline
        Informe final actual & \path{Informe_final/informe_final.pdf} & Activo & Evaluación académica principal. \\
        Manual de usuario & \path{Informe_final/Manual_de_usuario/manual_usuario.pdf} & Activo & Soporte al flujo visible. \\
        Manual técnico & \path{Informe_final/Manual_tecnico/manual_tecnico.pdf} & Activo & Soporte de arquitectura. \\
        Landing y docs & \path{docs/index.html}, \path{docs/manuals/} & Activo & Punto de acceso a build y docs. \\
        Validación e inst. & \path{Informe_final/validacion/} & Activo & Evidencia pública anonimizada. \\
        Telemetría agregada & \path{Telemetria/Mediciones_WebGL/} & Activo & Exportaciones profiler. \\
        Diagramas Mermaid & \path{Informe_final/figures/chapter4/} & Activo & Explican arquitectura y flujo. \\
        Anexos técnicos & \path{Informe_final/Desarrollo_App/Documentacion_Tecnica/} & Activo & Documentan WebGL y térmica. \\ \hline
    \end{tabular}
    \apanote{Esta tabla resume únicamente artefactos activos para consulta pública del cierre; auditorías históricas, bitácoras de trabajo, borradores y registros individuales se conservan como evidencia interna o por canal académico controlado, pero no como fuente primaria pública. Elaboración propia.}
\end{table}


\end{document}
